% ---------------------------------------------------------------------------
% Preambule pro přípravu na bakalářské státnice (MFF UK, Informatika / UI / SU)
% Sazba: pdflatex (UTF-8, čeština). Build: latexmk -pdf src/main.tex
% ---------------------------------------------------------------------------
\usepackage[T1]{fontenc}
\usepackage[utf8]{inputenc}
\usepackage{lmodern}
\usepackage[czech]{babel}
\usepackage[a4paper,margin=2.3cm]{geometry}

\usepackage{amsmath,amssymb,mathtools}
\usepackage{textcomp}
\usepackage{microtype}
\usepackage{xcolor}
\usepackage[export]{adjustbox}   % klíče max width / max height u \includegraphics
\usepackage{graphicx}
\graphicspath{{../}}             % obrázky z poznámek (relativně ke kořeni repa)
\usepackage{enumitem}
\usepackage{titlesec}
\usepackage[most]{tcolorbox}
\usepackage{booktabs}
\usepackage{multicol}
\usepackage{parskip}
\usepackage{hyperref}

% --- chybějící / nestandardní makra z poznámek (psány pro webový renderer) ---
\newcommand{\gt}{>}
\newcommand{\lt}{<}
\providecommand{\set}[1]{\{\,#1\,\}}
\renewcommand{\set}[1]{\{\,#1\,\}}
\providecommand{\braket}[1]{\left\langle #1 \right\rangle}
\renewcommand{\braket}[1]{\left\langle #1 \right\rangle}
\providecommand{\quotedblbase}{\glqq}
\providecommand{\quotesinglbase}{\glq}

% --- hlubší zanoření seznamů (poznámky jdou až do ~7 úrovní) ---
\setlistdepth{9}
\setlist{leftmargin=1.25em, itemsep=1.2pt, topsep=2.5pt, parsep=0pt}
\setlist[itemize,1]{label=\textbullet}
\setlist[itemize,2]{label=\normalfont\bfseries\textendash}
\setlist[itemize,3]{label=\textperiodcentered}
\setlist[itemize,4]{label=$\ast$}
\setlist[itemize,5]{label=$\cdot$}
\setlist[itemize,6]{label=$\circ$}
\setlist[itemize,7]{label=$\diamond$}
\renewlist{itemize}{itemize}{9}
\setlist[itemize,1]{label=\textbullet}
\setlist[itemize,2]{label=\normalfont\bfseries\textendash}
\setlist[itemize,3]{label=\textperiodcentered}
\setlist[itemize,4]{label=$\ast$}
\setlist[itemize,5]{label=$\cdot$}
\setlist[itemize,6]{label=$\circ$}
\setlist[itemize,7]{label=$\diamond$}
\setlist[itemize,8]{label=$\triangleright$}
\setlist[itemize,9]{label=$-$}

% --- barvy a styl nadpisů ---
\definecolor{mffblue}{RGB}{27,54,93}
\definecolor{mffaccent}{RGB}{0,90,156}
\definecolor{practiceyellow}{RGB}{180,120,0}
\definecolor{gapred}{RGB}{150,30,30}

\titleformat{\section}{\Large\bfseries\color{mffblue}}{\thesection}{0.6em}{}
\titleformat{\subsection}{\large\bfseries\color{mffaccent}}{\thesubsection}{0.6em}{}
\titleformat{\subsubsection}{\normalsize\bfseries}{\thesubsubsection}{0.6em}{}
\titlespacing*{\section}{0pt}{1.4em}{0.6em}

\hypersetup{
  colorlinks=true, linkcolor=mffaccent, urlcolor=mffaccent, citecolor=mffaccent,
  pdftitle={Příprava na bakalářské státnice -- Informatika / UI / Strojové učení},
  pdfauthor={Viktor Číhal}, bookmarksnumbered=true,
}

% --- callout boxy ---
% Procvičit: místo, kde poznámky odkazují na vlastní procvičení (není to mezera).
\newtcolorbox{practicebox}{
  enhanced, breakable, colback=practiceyellow!7, colframe=practiceyellow,
  boxrule=0.7pt, left=6pt, right=6pt, top=4pt, bottom=4pt, arc=2pt,
  before skip=5pt, after skip=5pt,
  overlay first={\node[anchor=north west,font=\bfseries\footnotesize\color{practiceyellow}]
    at ([xshift=6pt,yshift=-2pt]frame.north west){};},
  fonttitle=\bfseries\footnotesize, title={\textsc{Procvičit (poznámky odkazují na vlastní trénink)}}
}
% Mezera / doplněno: téma není v poznámkách (nebo jen jako alternativa).
\newtcolorbox{gapbox}{
  enhanced, breakable, colback=gapred!6, colframe=gapred,
  boxrule=0.8pt, left=6pt, right=6pt, top=4pt, bottom=4pt, arc=2pt,
  before skip=6pt, after skip=6pt,
  fonttitle=\bfseries\footnotesize, title={\textsc{Mezera v poznámkách -- ověř / doplň}}
}
% Info / shrnutí
\newtcolorbox{infobox}[1][]{
  enhanced, breakable, colback=mffaccent!5, colframe=mffaccent,
  boxrule=0.7pt, left=6pt, right=6pt, top=4pt, bottom=4pt, arc=2pt,
  before skip=6pt, after skip=6pt, fonttitle=\bfseries, #1
}

\setcounter{tocdepth}{2}
\setcounter{secnumdepth}{3}
\renewcommand{\arraystretch}{1.25}

% --- méně přetékajících řádků u dlouhých kódových úseků ---
\emergencystretch=3em
\makeatletter
\g@addto@macro\ttfamily{\hyphenchar\font=`\-}  % dovol dělení i ve strojopisu
\makeatother
\hbadness=10000  % netiskni varování o podtečení
